

\chapter{Results and Discussion}

In Chapter 5, students provide evidence of the "full functioning of working model as per stated requirements". This chapter is critical for the "Review 2" assessment, as it validates the "implementation and testing of the solution". 

\section{Working Prototype Demonstration}
Screenshots and descriptions of the functional model. This section serves as the visual proof of the "working model".

Interface Walkthrough: Include high-quality screenshots of the web application’s frontend, including input forms and result displays.

System Functionality: Provide a step-by-step description of how a user interacts with the ML model through the web interface.

Feature Verification: Map each screenshot to the initial "stated requirements" to prove the prototype is fully functional.

User Feedback (if any): Describe how the interface facilitates "societal impact" by being accessible to the target users.

\section{Validation and Inference}
Drawing proper inferences from experimental results or simulations. Students must demonstrate their ability to "draw the proper inferences from available results".

Experimental Results: Present data obtained through "theoretical, experimental, or simulations".

Justification of Results: Provide a "proper justification" for why the model produced specific outputs.

Inference Drawing: Explain what the results mean in the context of the "societal/research needs" identified in Chapter 1.

Comparative Validation: Briefly discuss how these results compare to the "proposed multiple solutions" considered during the design phase.


\section{Performance metrics} 

Accuracy, cost-effectiveness, and efficiency analysis. In this section, we explain the performance metrics used to check the performance of model. These metrices include accuracy, precision, recall, and F1-score. This section quantifies the "excellence" and "efficiency" of the engineering solution.

ML Accuracy: Use standard metrics such as Accuracy, Precision, Recall, and F1-Score to evaluate the "skill sets" applied in model training.

Confusion Matrix: Include a visual representation of the model's classification performance to support "validation with proper justification".

Cost-Effectiveness: Provide an analysis of "cost effectiveness" by comparing the performance gained against the resources/computation used.

System Efficiency: Analyze the "impact of solutions" in terms of response time and resource efficiency for "sustainable development".